%% Expert Systems with Applications --- Elsevier CAS single-column template
\documentclass[a4paper,fleqn]{cas-sc}

\usepackage[authoryear,longnamesfirst]{natbib}
\usepackage{amsmath,amssymb}

\begin{document}
\let\WriteBookmarks\relax
\def\floatpagepagefraction{1}
\def\textpagefraction{.001}

\shorttitle{Logarithmic Weighted Cross-Entropy for Imbalanced Classification}
\shortauthors{S.~H.~Go et~al.}

\title[mode=title]{Logarithmic Weighted Cross-Entropy: A Stable and
Parameter-Efficient Re-weighting Loss for Class-Imbalanced Classification}

\author[1]{Seung Ho Go}
\ead{gseungho999@gmail.com}
\credit{Conceptualization, Methodology, Software, Writing -- original draft}

\affiliation[1]{organization={Sejong University},
            addressline={209 Neungdong-ro, Gwangjin-gu},
            city={Seoul},
            country={Republic of Korea}}

\author[1]{Hyunseok Lee}
\credit{Methodology, Validation, Writing -- review \& editing}

\author[1]{Jin Hee Yoon}
\cormark[1]
\credit{Supervision, Conceptualization, Writing -- review \& editing}

\cortext[1]{Corresponding author}

\begin{abstract}
Class imbalance is a pervasive obstacle in machine learning, arising whenever a
few majority categories dominate the training data while critically important
minority categories remain severely underrepresented. Among algorithm-level
remedies, Weighted Cross-Entropy (WCE) is the most direct: it scales the loss of
each class by the inverse of its frequency. However, these inverse-frequency
weights grow without bound as the imbalance becomes more severe, producing
unstable optimization. More elaborate alternatives such as Class-Balanced (CB)
Loss and Focal Loss introduce additional hyperparameters and, in our experiments,
behave inconsistently across datasets---occasionally falling below ordinary
Cross-Entropy (CE). We propose Logarithmic Weighted Cross-Entropy (LWCE), which
replaces the inverse-frequency weight with a logarithmically compressed
counterpart that keeps a meaningful boost for rare classes while bounding weight
growth, and which introduces no additional hyperparameters. We further generalize
it to Power-Log Weighted Cross-Entropy (PLWCE) through a single exponent that
adjusts the focusing strength and is selected automatically on a small proxy
subset. To assess generality, we evaluate the proposed losses across three
distinct application domains that together span a wide range of imbalance
severity: long-tailed image classification, tabular classification benchmarks,
and network intrusion detection. Across these domains, LWCE and PLWCE improve
over Cross-Entropy as the imbalance grows and, unlike the competing re-weighting
losses, never collapse to substantially worse performance on any dataset. These
results indicate that logarithmic weight compression offers a simple, stable, and
parameter-efficient alternative to existing re-weighting strategies for
class-imbalanced classification.
\end{abstract}

\begin{highlights}
\item A logarithmic re-weighting loss that bounds the weight explosion of weighted cross-entropy.
\item A single-parameter generalization tuned automatically on a small proxy subset.
\item Evaluation spans image, tabular, and network-intrusion domains of varying imbalance.
\item The proposed losses improve as imbalance grows and never collapse on any dataset.
\end{highlights}

\begin{keywords}
Class imbalance \sep Long-tailed learning \sep Cross-entropy loss \sep
Cost-sensitive learning \sep Network intrusion detection
\end{keywords}

\maketitle

\section{Introduction}\label{sec:intro}

Across a wide range of real-world classification tasks---including medical
diagnosis, financial fraud detection, network intrusion detection, and visual
recognition---the categories of interest are rarely observed in equal
proportions. Instead, a small number of majority categories concentrate the
overwhelming majority of the available samples, while many minority categories
are represented by only a handful of examples~\citep{he2009,johnson2019}. This
long-tailed structure is not an artifact of poor data collection but an intrinsic
property of the underlying phenomena: healthy patients vastly outnumber those
with a rare disease, legitimate transactions dwarf fraudulent ones, and benign
network traffic far exceeds malicious traffic. Because the minority categories
are often precisely the ones that carry the greatest practical cost when missed,
building classifiers that remain reliable under imbalance is a problem of broad
and lasting importance.

The central difficulty is that the standard training objective is itself biased
toward the majority. When a model is trained with ordinary Cross-Entropy (CE)
loss, the aggregate gradient signal is dominated by the numerically abundant
classes, so the optimizer can drive the overall loss down while largely ignoring
rare categories. The resulting classifier attains a high overall accuracy yet
performs poorly on exactly the minority classes that matter most---a failure mode
commonly described as majority bias. Standard accuracy therefore masks the
problem rather than revealing it, and a method that addresses imbalance must
improve minority-class performance without sacrificing the majority.

Existing remedies fall broadly into two families. Resampling methods reshape the
training distribution directly: oversampling techniques such as the Synthetic
Minority Over-sampling Technique (SMOTE) synthesize new minority
instances~\citep{chawla2002}, while undersampling discards majority examples to
restore balance. Both directions carry well-known risks---oversampling can
introduce synthetic artifacts and encourage overfitting, whereas undersampling
discards potentially informative majority data~\citep{drummond2003}. Algorithm-level
methods instead modify the training objective. These include margin-based losses
that enlarge the decision boundary around minority
classes~\citep{cao2019}, logit-adjustment techniques that calibrate class priors,
and loss re-weighting methods that scale the per-class loss by class-dependent
weights. Among these, re-weighting is especially attractive because it is
architecture-agnostic, modality-independent, and requires only a minimal change
to the loss function; this paper focuses on the re-weighting family.

Despite its appeal, re-weighting remains unsatisfying in practice. Weighted
Cross-Entropy (WCE), the most direct re-weighting scheme, assigns each class a
weight inversely proportional to its frequency; as the imbalance grows more
severe these weights grow without bound, producing extreme gradients that
destabilize training. More sophisticated alternatives were designed to soften
this behavior but bring their own drawbacks. Focal Loss reduces the influence of
easy, typically majority, examples through a tunable focusing term, but it
operates on per-sample confidence rather than class frequency and therefore
behaves inconsistently across imbalance levels. Class-Balanced (CB) Loss
re-weights by an effective number of samples governed by an additional
hyperparameter, yet that hyperparameter is sensitive and, in our experiments, CB
Loss degrades below plain Cross-Entropy under severe imbalance. In short, the
field still lacks a re-weighting loss that is simultaneously simple,
hyperparameter-light, numerically stable, and dependable across both the level of
imbalance and the application domain.

To address this gap, we propose Logarithmic Weighted Cross-Entropy (LWCE). The
key idea is to replace the inverse-frequency weight of WCE with a logarithmically
compressed counterpart. The logarithm fundamentally bounds the growth of the
weights, so the explosive behavior of WCE is removed while a meaningful emphasis
on rare classes is preserved; crucially, the resulting loss introduces no
additional hyperparameters and can be used as a drop-in replacement for
Cross-Entropy. We further generalize LWCE to Power-Log Weighted Cross-Entropy
(PLWCE) by introducing a single exponent that continuously controls the focusing
strength, ranging from near-uniform weighting to an aggressive emphasis on the
tail. This exponent is the only free parameter, and it is selected automatically
through a lightweight grid search on a small proxy subset, adding negligible
computational overhead.

We assess the proposed losses through a deliberately broad empirical study rather
than a single benchmark. To test whether the benefit of logarithmic compression
generalizes, we evaluate across three application domains that together cover a
wide spectrum of imbalance severity: long-tailed image classification on standard
vision benchmarks, a collection of tabular classification datasets drawn from
public repositories, and several network intrusion detection datasets in which
rare attack types are extremely scarce. In every domain we compare the same set
of re-weighting losses under a controlled protocol with repeated runs, and we
report metrics that explicitly reflect minority-class performance rather than
overall accuracy alone. This cross-domain design lets us examine not only whether
the proposed losses help, but how their advantage changes as imbalance becomes
more extreme.

The main contributions of this paper are as follows.
\begin{enumerate}
  \item We identify and characterize the weight-explosion failure mode of Weighted
        Cross-Entropy and motivate logarithmic compression as a principled remedy.
  \item We propose Logarithmic Weighted Cross-Entropy, a parameter-free drop-in
        replacement for Weighted Cross-Entropy, together with its single-parameter
        generalization, Power-Log Weighted Cross-Entropy, whose only exponent is
        tuned automatically on a small proxy subset.
  \item We provide a controlled cross-domain evaluation spanning image, tabular,
        and network intrusion detection data, showing that the proposed losses
        improve over Cross-Entropy as imbalance grows and, unlike Class-Balanced
        Loss and Focal Loss, never collapse to substantially worse performance on
        any dataset.
\end{enumerate}

The remainder of this paper is organized as follows.
Section~\ref{sec:related} reviews related work on learning from imbalanced and
long-tailed data. Section~\ref{sec:method} defines the proposed losses and
analyzes their properties. Section~\ref{sec:experiments} describes the
experimental setup across the three domains and presents the results.
Section~\ref{sec:conclusion} concludes and discusses limitations and future work.

\section{Related Work}\label{sec:related}

\section{Proposed Method}\label{sec:method}

\section{Experiments}\label{sec:experiments}

\section{Conclusion}\label{sec:conclusion}

\printcredits

\bibliographystyle{cas-model2-names}
\bibliography{refs}

\end{document}
